\begin{proof}[Proof of \cref{obj:lem:walsh-channel-energy}]
\begin{enumerate}
\item For every subset \(E\subseteq\{1,\ldots,d\}\), define its sign character by
\[
\psi_E(\varepsilon)=\prod_{j\in E}\varepsilon_j,
\qquad \varepsilon\in\{-1,1\}^d.
\]
Also define the initial coordinate sets
\[
E_s=\{1,\ldots,\min\{s,d\}\},
\qquad s\in\mathbb Z_{\ge0}.
\]
Fix \(w\in\mathbb R\) with \(g_d(w)>0\). The empty-product convention and the definition of \(g_d\) give
\[
a_0(w)
=2^{-d}\sum_{\varepsilon\in\{-1,1\}^d}
\frac{F_\varepsilon(w)}{g_d(w)}
=\frac{g_d(w)}{g_d(w)}
=1.
\]

Relabeling coordinates preserves the sign sum and hence \(F_\varepsilon(w)\). More explicitly, for any permutation
\[
\rho:\{1,\ldots,d\}\longrightarrow\{1,\ldots,d\},
\qquad
(\varepsilon\circ\rho)_j=\varepsilon_{\rho(j)},
\]
one has
\[
F_{\varepsilon\circ\rho}(w)=F_\varepsilon(w),
\qquad
\psi_{\rho(E)}(\varepsilon)=\psi_E(\varepsilon\circ\rho).
\]
For a given \(E\), choose such a permutation with
\(\rho(E_{|E|})=E\), which is possible because both sets have cardinality \(|E|\). Reindexing the sign sum therefore yields
\[
a_{|E|}(w)
=2^{-d}\sum_{\varepsilon\in\{-1,1\}^d}
\frac{F_\varepsilon(w)}{g_d(w)}\psi_E(\varepsilon).
\]

The finite character kernel satisfies
\[
\sum_{E\subseteq\{1,\ldots,d\}}
\psi_E(\delta)\psi_E(\varepsilon)
=
\prod_{j=1}^d(1+\delta_j\varepsilon_j)
=
\begin{cases}
2^d,&\delta=\varepsilon,\\
0,&\delta\ne\varepsilon,
\end{cases}
\qquad \varepsilon,\delta\in\{-1,1\}^d.
\]
The first equality expands the product over all subsets. If the sign vectors agree, every factor equals two; otherwise a coordinate where they differ supplies a zero factor. Substituting the coefficient formula and using this kernel gives
\[
\begin{aligned}
\sum_{E\subseteq\{1,\ldots,d\}}
a_{|E|}(w)\psi_E(\varepsilon)
&=
2^{-d}\sum_{\delta\in\{-1,1\}^d}
\frac{F_\delta(w)}{g_d(w)}
\sum_{E\subseteq\{1,\ldots,d\}}
\psi_E(\delta)\psi_E(\varepsilon)\\
&=\frac{F_\varepsilon(w)}{g_d(w)}.
\end{aligned}
\]
% lean: CausalSmith.Experimentation.ThinnedgraphAdditiveRiskFrontier.rowDensity_walsh_expansion

\item Each \(F_\varepsilon\) is a translate of the baseline probability density, so
\[
F_\varepsilon(w)\ge0,
\qquad
\int_{\mathbb R}F_\varepsilon(w)\,dw=1.
\]
Consequently \(g_d(w)\ge0\), and \(g_d(w)=0\) forces every \(F_\varepsilon(w)=0\). Thus cancellation of the denominator, including the zero-density convention, gives the identity
\[
a_s(w)g_d(w)
=
2^{-d}\sum_{\varepsilon\in\{-1,1\}^d}
F_\varepsilon(w)\psi_{E_s}(\varepsilon),
\qquad w\in\mathbb R,\quad s\in\mathbb Z_{\ge0}.
\]
For \(s\ge1\), the assumptions \(d\ge1\) and \(s\ge1\) ensure that \(E_s\) contains coordinate one, including when \(s>d\). Summing signs coordinate by coordinate gives
\[
\sum_{\varepsilon\in\{-1,1\}^d}\psi_{E_s}(\varepsilon)
=
\prod_{j=1}^d
\left[
\sum_{\varepsilon_j\in\{-1,1\}}
\begin{cases}
\varepsilon_j,&j\in E_s,\\
1,&j\notin E_s
\end{cases}
\right]
=0,
\]
because the factor for coordinate one is \((-1)+1=0\). Integrating the preceding cancellation identity, with the finite sum exchanged with the integral, now yields
\[
\int_{\mathbb R}a_s(w)g_d(w)\,dw
=
2^{-d}\sum_{\varepsilon\in\{-1,1\}^d}
\psi_{E_s}(\varepsilon)
\int_{\mathbb R}F_\varepsilon(w)\,dw
=0.
\]
% lean: CausalSmith.Experimentation.ThinnedgraphAdditiveRiskFrontier.walshCoeff_centered

\item Nonnegativity of \(g_d\) gives
\[
\gamma_s=\int_{\mathbb R}a_s(w)^2g_d(w)\,dw\ge0
\qquad(1\le s\le d).
\]
Since \(s\ge1\) on this range,
\[
0\le \binom ds\gamma_s
\le s\binom ds\gamma_s.
\]
Summing proves
\[
0\le\eta\le\eta_1.
\]
% lean: CausalSmith.Experimentation.ThinnedgraphAdditiveRiskFrontier.eta_nonneg_le_etaOne

\item First consider \(h=0\). Then \(F_\varepsilon(w)=f(w)\) for every sign vector. At a point where \(g_d(w)=0\), the coefficient convention gives \(a_s(w)=0\). At any other point, for every \(s\ge1\), the character cancellation from the preceding argument gives
\[
a_s(w)
=
2^{-d}\frac{f(w)}{g_d(w)}
\sum_{\varepsilon\in\{-1,1\}^d}\psi_{E_s}(\varepsilon)
=0.
\]
Hence every \(\gamma_s\), \(1\le s\le d\), vanishes, and
\[
\eta_1=0=\frac{12\pi^2h^2}{d}.
\]
% lean: CausalSmith.Experimentation.ThinnedgraphAdditiveRiskFrontier.etaOne_zero_amplitude

\item Suppose henceforth that \(h\ne0\). Define the coordinate reversal and its averaged density energy by
\[
(\varepsilon^{(j)})_i=
\begin{cases}
-\varepsilon_j,&i=j,\\
\varepsilon_i,&i\ne j,
\end{cases}
\qquad
\varepsilon\in\{-1,1\}^d,\quad i,j\in\{1,\ldots,d\},
\]
and
\[
\mathcal D_j(w)=
\begin{cases}
2^{-d}\displaystyle\sum_{\varepsilon\in\{-1,1\}^d}
\frac{(F_\varepsilon(w)-F_{\varepsilon^{(j)}}(w))^2}{g_d(w)},
&g_d(w)>0,\\
0,&g_d(w)=0,
\end{cases}
\qquad w\in\mathbb R,\quad 1\le j\le d.
\]

To derive the weighted identity, let
\[
v_{\mathrm{cube}}:\{-1,1\}^d\longrightarrow\mathbb R,
\qquad
\widehat v_{\mathrm{cube}}(E)
=
2^{-d}\sum_{\varepsilon\in\{-1,1\}^d}
v_{\mathrm{cube}}(\varepsilon)\psi_E(\varepsilon),
\quad E\subseteq\{1,\ldots,d\}.
\]
Expanding the squared coefficients and using the character kernel established above proves finite Parseval:
\[
\begin{aligned}
\sum_{E\subseteq\{1,\ldots,d\}}
\widehat v_{\mathrm{cube}}(E)^2
&=
2^{-2d}\sum_{\varepsilon,\delta\in\{-1,1\}^d}
v_{\mathrm{cube}}(\varepsilon)v_{\mathrm{cube}}(\delta)
\sum_{E\subseteq\{1,\ldots,d\}}
\psi_E(\varepsilon)\psi_E(\delta)\\
&=
2^{-d}\sum_{\varepsilon\in\{-1,1\}^d}
v_{\mathrm{cube}}(\varepsilon)^2.
\end{aligned}
\]
A coordinate reversal is an involution, and
\[
\psi_E(\varepsilon^{(j)})
=
\begin{cases}
-\psi_E(\varepsilon),&j\in E,\\
\psi_E(\varepsilon),&j\notin E.
\end{cases}
\]
Reindexing by this involution therefore gives
\[
2^{-d}\sum_{\varepsilon\in\{-1,1\}^d}
\bigl(v_{\mathrm{cube}}(\varepsilon)
-v_{\mathrm{cube}}(\varepsilon^{(j)})\bigr)\psi_E(\varepsilon)
=
\begin{cases}
2\widehat v_{\mathrm{cube}}(E),&j\in E,\\
0,&j\notin E.
\end{cases}
\]
Applying finite Parseval to the coordinate difference yields
\[
\sum_{\substack{E\subseteq\{1,\ldots,d\}\\j\in E}}
\widehat v_{\mathrm{cube}}(E)^2
=
\frac14\,2^{-d}\sum_{\varepsilon\in\{-1,1\}^d}
\bigl(v_{\mathrm{cube}}(\varepsilon)
-v_{\mathrm{cube}}(\varepsilon^{(j)})\bigr)^2.
\]
Summing over \(j\) counts each subset exactly \(|E|\) times, so
\[
\sum_{E\subseteq\{1,\ldots,d\}}
|E|\widehat v_{\mathrm{cube}}(E)^2
=
\frac14\sum_{j=1}^d2^{-d}
\sum_{\varepsilon\in\{-1,1\}^d}
\bigl(v_{\mathrm{cube}}(\varepsilon)
-v_{\mathrm{cube}}(\varepsilon^{(j)})\bigr)^2.
\]

For \(g_d(w)>0\), apply this identity with
\[
v_{\mathrm{cube}}(\varepsilon)=\frac{F_\varepsilon(w)}{g_d(w)}.
\]
Its subset coefficients are \(a_{|E|}(w)\). Grouping subsets by cardinality and multiplying by \(g_d(w)\) gives
\[
\sum_{s=0}^d s\binom ds a_s(w)^2g_d(w)
=
\frac14\sum_{j=1}^d\mathcal D_j(w).
\]
At \(g_d(w)=0\), both sides vanish by their conventions, so this identity holds for every \(w\).

The coefficient densities are integrable: the triangle inequality gives, at positive \(g_d(w)\),
\[
|a_s(w)|
\le
2^{-d}\sum_{\varepsilon\in\{-1,1\}^d}
\frac{F_\varepsilon(w)}{g_d(w)}
=1,
\]
and the zero-density convention preserves this bound. Thus
\[
0\le a_s(w)^2g_d(w)\le g_d(w),
\qquad
\int_{\mathbb R}g_d(w)\,dw=1.
\]
The flip energies are integrable by the bounded compact-support envelope established in the next step. Integrating the pointwise identity and discarding its zero \(s=0\) term therefore gives
\[
\eta_1=\frac14\sum_{j=1}^d\int_{\mathbb R}\mathcal D_j(w)\,dw.
\]
% lean: CausalSmith.Experimentation.ThinnedgraphAdditiveRiskFrontier.walsh_weighted_integral_parseval

\item By \cref{obj:lem:baseline-translation-affinity}, the square root of the baseline density is globally \(4\pi\)-Lipschitz. Hence, for \(w_{\mathrm a},w_{\mathrm b}\in\mathbb R\),
\[
\bigl(\sqrt{f(w_{\mathrm a})}-\sqrt{f(w_{\mathrm b})}\bigr)^2
\le16\pi^2(w_{\mathrm a}-w_{\mathrm b})^2.
\]
Also,
\[
\bigl(\sqrt{f(w_{\mathrm a})}+\sqrt{f(w_{\mathrm b})}\bigr)^2
\le2\bigl(f(w_{\mathrm a})+f(w_{\mathrm b})\bigr),
\]
as follows by expanding the nonnegative square of their difference. Factoring the difference of the densities and multiplying these bounds gives
\[
\begin{aligned}
\bigl(f(w_{\mathrm a})-f(w_{\mathrm b})\bigr)^2
&=
\bigl(\sqrt{f(w_{\mathrm a})}-\sqrt{f(w_{\mathrm b})}\bigr)^2
\bigl(\sqrt{f(w_{\mathrm a})}+\sqrt{f(w_{\mathrm b})}\bigr)^2\\
&\le
32\pi^2(w_{\mathrm a}-w_{\mathrm b})^2
\bigl(f(w_{\mathrm a})+f(w_{\mathrm b})\bigr).
\end{aligned}
\]

Define the translation amounts by
\[
\zeta_\varepsilon=\frac{h}{2d}\sum_{i=1}^d\varepsilon_i,
\qquad \varepsilon\in\{-1,1\}^d.
\]
Reversing coordinate \(j\) gives
\[
\sum_{i=1}^d(\varepsilon^{(j)})_i
=\sum_{i=1}^d\varepsilon_i-2\varepsilon_j,
\qquad
\zeta_{\varepsilon^{(j)}}=\zeta_\varepsilon-\frac hd\varepsilon_j.
\]
Consequently the two translated arguments have squared difference \(h^2/d^2\), and the preceding density bound becomes
\[
\bigl(F_\varepsilon(w)-F_{\varepsilon^{(j)}}(w)\bigr)^2
\le
\frac{32\pi^2h^2}{d^2}
\bigl(F_\varepsilon(w)+F_{\varepsilon^{(j)}}(w)\bigr).
\]
Since the reversal permutes the sign cube,
\[
2^{-d}\sum_{\varepsilon\in\{-1,1\}^d}
\bigl(F_\varepsilon(w)+F_{\varepsilon^{(j)}}(w)\bigr)
=2g_d(w).
\]
Summing the density bound and dividing by \(g_d(w)>0\) proves
\[
0\le\mathcal D_j(w)\le\frac{64\pi^2h^2}{d^2}.
\]
At \(g_d(w)=0\), the same inequality follows from \(\mathcal D_j(w)=0\).

For the support bound, define
\[
I_h=\left[-\frac14-\frac h2,\frac14+\frac h2\right].
\]
The triangle inequality gives
\[
|\zeta_\varepsilon|
\le\frac{h}{2d}\sum_{i=1}^d|\varepsilon_i|
=\frac h2.
\]
Thus, if \(w\notin I_h\),
\[
|w-\zeta_\varepsilon|
\ge |w|-|\zeta_\varepsilon|>\frac14,
\]
so the support of \(f\) specified in
\cref{obj:lem:baseline-translation-affinity} implies
\(F_\varepsilon(w)=0\) for every \(\varepsilon\), and hence \(g_d(w)=0\).
Writing the interval indicator explicitly as
\[
\mathbf1_{I_h}(w)=
\begin{cases}
1,&w\in I_h,\\
0,&w\notin I_h,
\end{cases}
\]
we obtain the integrable envelope
\[
0\le\mathcal D_j(w)
\le\frac{64\pi^2h^2}{d^2}\mathbf1_{I_h}(w).
\]
The interval has length \(1/2+h\le3/4\). Therefore
\[
\int_{\mathbb R}\mathcal D_j(w)\,dw
\le
\left(\frac12+h\right)\frac{64\pi^2h^2}{d^2}
\le\frac{48\pi^2h^2}{d^2}.
\]
Finally, summing this bound over the \(d\) coordinates in the integrated identity gives
\[
\eta_1
=\frac14\sum_{j=1}^d\int_{\mathbb R}\mathcal D_j(w)\,dw
\le\frac d4\,\frac{48\pi^2h^2}{d^2}
=\frac{12\pi^2h^2}{d}.
\]
Together with the preceding steps, this proves all the assertions.
% lean: CausalSmith.Experimentation.ThinnedgraphAdditiveRiskFrontier.integrated_flip_sum_le
\end{enumerate}
\end{proof}
